\section{Filtrations, localisation and group actions} \label{section:localisation-filtration}

\subsection{Recollections on filtered vector spaces and algebras}

An \emph{increasing filtration} on a vector space $V$ is the data of subspaces $V_p$ of $V$, for all $p$ in $\Z$, such that $V_p \subseteq V_{p+1}$.

A filtration is called:
\begin{itemize}
	\item \emph{separated} if $0 = \bigcap_{p \in \Z} V_p$,
	
	\item \emph{bounded from below} if $V_p = 0$ for some $p$ in $\Z$,
	
	\item \emph{bounded from above} if $V_p = V$ for some $p$ in $\Z$,
	
	\item \emph{exhaustive} if $V = \bigcup_{p \in \Z} V_p$.
\end{itemize}

The associated graded vector space of a filtered vector space $V$ is 
\[ \gr V \coloneqq \bigoplus_{p \in \Z} V_p / V_{p - 1}. \]
Let $V$ and $W$ be two filtered vector spaces. A linear map $\phi \colon V \rightarrow W$ is called a \emph{filtered homomorphism} if $f(V_p) \subseteq W_p$ for all $p$ in $\Z$. Then we can define an associated graded homomorphism
\[ \gr \phi \colon \gr V \longrightarrow \gr W \]
in an obvious way.

The following fact is standard, see \cite[Lemma 6.8]{kamnitzer2022hamiltonian} for a proof.

\begin{lemma} \label{lemma:quantisation-iso}
	Let $V$ and $W$ be two filtered vector spaces such that
	\begin{enumerate}
		\item both filtrations are exhaustive
		
		\item the filtration on $V$ is separated,
		
		\item the filtration on $W$ is bounded from below.
	\end{enumerate} 
	Let $\phi \colon V \rightarrow W$ be a filtered homomorphism. Therefore, if $\gr \phi \colon \gr V \rightarrow \gr W$ is an isomorphism, then $\phi \colon V \rightarrow W$ is an isomorphism.
\end{lemma}

The assignment
\[ (V \otimes W)_p \coloneqq \sum_{i + j = p} V_i \otimes W_j, \quad \text{for} \quad p \in \Z, \]
defines a filtration on the tensor product $V \otimes W$. There is a natural isomorphism
\[ \gr V \otimes \gr W \cong \gr(V \otimes W). \]

\subsection{Localisation of a filtered algebra} \label{subsectin:localisation-filtered}


In this section, we recall how one can localise a positively filtered algebra according to \cite[Section 2.1]{kapranov1998noncommutative}. See also \cite[Section 6]{kamnitzer2022hamiltonian} or \cite{skoda2006noncommutative}.


Let $A$ be a unital associative algebra equipped with an increasing filtration by subspaces $A_p$ ($p$ is in $\Z$), such that $A_p = 0$ for $p < 0$, and we have the following multiplicative properties:
\begin{equation} \label{equation:multiplicative-properties} 1 \in A_0, \quad A_p \cdot A_q \subseteq A_{p + q} \quad \text{and} \quad [A_p, A_q] \subseteq A_{p + q - 1}, \end{equation}
for $p$ and $q$ in $\Z$. One sees that $A_0$ is a commutative subalgebra and the associated graded algebra
\[ \gr A \coloneqq \bigoplus_{p \in \Z} A_p / A_{p - 1} = \bigoplus_{p \geqslant 0} A_p / A_{p - 1} \]
is commutative. Moreover, $A_0$ identifies with the degree $0$ subalgebra of $\gr A$.  %Moreover, $A$ is equipped with two commuting left and right $A_0$-actions respectively by left and right multiplication. Let $e$ be in $A_0$ and denote by 
%\[ A[e^{-1}] \coloneqq A_0[e^{-1}] \otimes_{A_0} A \] 
%the left localisation with respect to $e$, that is a left $A_0[e^{-1}]$ and a right $A_0$-module.

The following lemma and its proof is inspired by \cite[Proposition 2.1.5]{kapranov1998noncommutative}

\begin{lemma} \label{lemma:ore}
	Any multiplicative subset $S$ in $A_0$ satisfies the left and right Ore condition in $A$:
	\begin{enumerate}
		\item \label{lemma:ore,item:left-frac} for any $a$ in $A$ and $s$ in $S$, there are $b$ in $A$ and $t$ in $S$ such that $t a = b s$,
		
		\item \label{lemma:ore,item:left-zero} if $a$ in $A$ and $s$ satisfy $a s = 0$, then there is some $t$ in $S$ such that $t a = 0$,
		
		\item for any $a$ in $A$ and $s$ in $S$, there are $b$ in $A$ and $t$ in $S$ such that $a t = s b$,
		
		\item if $a$ in $A$ and $s$ satisfy $s a = 0$, then there is some $t$ in $S$ such that $a t = 0$.
	\end{enumerate}
\end{lemma}

\begin{proof}
	We check only the left conditions \eqref{lemma:ore,item:left-frac} and \eqref{lemma:ore,item:left-zero}.
	
	Assume that $a$ is in $A_n$ for some nonnegative integer $n$ and $s$ is in $S$. Then for any positive integer $k$:
	\[ s^{k} a = \bigg(\sum_{i=0}^{k - 1} s^{k - 1 - i} \ad(s)^i a\bigg) s + \ad(s)^{k} a. \] 
	Note that $\ad(s)^{k} a$ belongs to $A_{n - k}$. So for $k = n + 1$,
	$\ad(s)^{n + 1} a = 0$ since $A_{-1} = 0$, and we get
	\[ s^{n+1} a = \bigg(\sum_{i=0}^{n} s^{n - i} \ad(s)^i a\bigg) s. \]
	Then \eqref{lemma:ore,item:left-frac} holds for
	\[ b \coloneqq \sum_{i=0}^{n} s^{n - i} \ad(s)^i a \quad \text{and} \quad t \coloneqq s^{n + 1}. \]
	
	Now assume that $a s = 0$. Then for any positive integer $k$:
	\[ s^k a = \ad(s)^k a. \]
	For $k = n + 1$, we get
	\[ s^{n + 1} a = \ad(s)^{n + 1} a = 0. \]
	Therefore \eqref{lemma:ore,item:left-zero} holds for $t = s^{n + 1}$.
\end{proof}

By the general theory of Ore localisation \cite[Chapter II]{stenstrom1975rings}, we get the following.

\begin{proposition} \label{proposition:fraction}
	For any multiplicative subset $S$ in $A_0$, there is an algebra $A[S^{-1}]$ equipped with an algebra homomorphism $\iota \colon A \rightarrow A[S^{-1}]$ satisfying the following universal property. For any algebra homomorphism $\psi \colon A \rightarrow B$ such that $\psi(S)$ lies in the group of units $B^\times$ of $B$, there is a unique algebra homomorphism $\widetilde{\psi} \colon A[S^{-1}] \rightarrow B$ such that $\widetilde{\psi} \circ \iota = \psi$.
	
	The algebra $A[S^{-1}]$ is a flat left and right $A$-module. The natural map
	\[ A_0[S^{-1}] \otimes_{A_0} A \longrightarrow A[S^{-1}] \]
	is an isomorphism of left $A_0[S^{-1}]$-module.
\end{proposition}

There is a natural filtration on $A[S^{-1}]$: because $A_0 \rightarrow A_0[S^{-1}]$ is flat, the chain of inclusions
\[ A_0 \subseteq A_1 \subseteq \cdots \subseteq A_p \subseteq \cdots \subseteq A \]
gives rise to a chain of inclusion
\[ A_0[S^{-1}] \subseteq A_0[S^{-1}] \otimes_{A_0} A_1  \subseteq \cdots \subseteq A_0[S^{-1}] \otimes_{A_0} A_p \subseteq \cdots \subseteq A[S^{-1}] \]
after tensoring by $A_0[S^{-1}]$ and using the identification $A_0[S^{-1}] \otimes_{A_0} A \cong A[S^{-1}]$. Therefore, the assignment
\[ (A[S^{-1}])_p \coloneqq A_0[S^{-1}] \otimes_{A_0} A_p, \quad \text{for} \quad p \in \Z,  \]
defines an increasing filtration on $A[S^{-1}]$ by vector subspaces. This filtration clearly satisfies the multiplicative properties \eqref{equation:multiplicative-properties}. In particular, $\gr(A[S^{-1}])$ is a graded commutative algebra and contains $A_0[S^{-1}]$ as the degree 0 subalgebra.

\begin{proposition} \label{proposition:gr-localisation}
	There is a natural algebra isomorphism
	\[ (\gr A)[S^{-1}] \longrightarrow \gr(A[S^{-1}]). \]
\end{proposition}

\begin{proof}
	The natural homomorphism $A \rightarrow A[S^{-1}]$ is compatible with the filtrations on both sides, therefore we get an algebra homomorphism
	\[ \gr A \longrightarrow \gr(A[S^{-1}]). \]
	Under this last homomorphism, $S$ on the left-hand side is mapped to the degree $0$ component of the right-hand side, that is to say $A_0[S^{-1}]$. Therefore, by universal property we get the natural algebra homomorphism
	\begin{equation} \label{equation:graded-iso} (\gr A)[S^{-1}] \longrightarrow \gr(A[S^{-1}]). \end{equation}
	
	To show that it is an isomorphism, we note that, for any $p$ in $\Z$, the exact sequence
	\[ 0 \longrightarrow A_{p - 1} \longrightarrow A_p \longrightarrow \gr_p A \longrightarrow 0 \]
	gives rise to an exact sequence
	\[ 0 \longrightarrow (A[S^{-1}])_{p - 1} \longrightarrow (A[S^{-1}])_p \longrightarrow A_0[S^{-1}] \otimes_{A_0} (\gr_p A) \longrightarrow 0 \]
	after applying $A_0[S^{-1}] \otimes_{A_0} \bullet$. Therefore, we get an isomorphism
	\[ \gr_p(A[S^{-1}]) \cong A_0[S^{-1}] \otimes_{A_0} (\gr_p A) \]
	of $A_0[S^{-1}]$-modules. This isomorphism gives an inverse to the natural algebra homomorphism \eqref{equation:graded-iso}, that is then an isomorphism.
\end{proof}

%\begin{lemma}
%	The right multiplication by $e$ gives an invertible endomorphism of the left localisation $A[e^{-1}]$. Therefore, $A[e^{-1}]$ is naturally a left and right $A_0[e^{-1}]$-module.
%\end{lemma}
%
%
%\begin{proof}
%	The right multiplication by $e$ is given by the formula:
%	\[  x e = e x + [x, e], \quad \text{for} \quad x \in A. \]
%	It gives a left $A_0$-linear endomorphism of $A$, that extends to a $A_0[e^{-1}]$-linear automorphism of $A[e^{-1}]$. Moreover, $[\bullet, e]$ is nilpotent because the brackets goes down the filtration. As a sum of an invertible and a nilpotent endomorphism, the right multiplication by $e$ induces an invertible endomorphism of $A[e^{-1}]$. This proves the statement.
%\end{proof}
%
%\begin{proposition}
%	There is a unique algebra structure on $A[e^{-1}]$ such that the natural homomorphism
%	\[ A \longrightarrow A[e^{-1}] \]
%	is an algebra homomorphism.
%\end{proposition}
%
%\begin{proof}
%	By localisation, the multiplication on $A$ gives rise to a map
%	
%	
%	It is clearly bilinear and associative.
%\end{proof}

\subsection{Filtered representations of an algebraic group}

\label{subsection:filtered-representation}

In this section, we recall from \cite{juillard2025phd} the notion of a filtered module over a graded algebraic group, and some useful properties to study the subspace of invariants. See also \cite[Section 2.4]{hamilton2026relative} and \cite[Section 3.1]{hoskins2025moduli}, where very similar notions are introduced.

Let $N$ be an algebraic group and $\n$ be its Lie algebra. The data of a Lie algebra $\Z$-grading on $\n$ is equivalent to the data of an algebraic action of $\Gm$ on $\n$ by Lie algebra automorphisms. By analogy, we define a \emph{grading on $N$} as the data of an action of $\Gm$ on $N$ by algebraic group automorphisms.


If $X$ is an affine scheme equipped with an action $\sigma \colon \Gm \times X \rightarrow X$, then $\C[X]$ is graded by
$$ \C[X](p) = \{F \in \C[X] ~ | ~  \sigma^*(F) = t^p \otimes F \} ~ \quad ~ \text{for} \quad p \in \Z. $$
Therefore, if $N$ is a graded group, $\C[N]$ is a graded algebra.

Let $N$ be a graded group and $V$ be an $N$-module, that is to say there is a co-action $V \rightarrow \C[N] \otimes V$ \cite[Chapter 4, Section a]{milne2017algebraic}. We say that an $N$-module $V$ is a \emph{graded $N$-module} if $V = \bigoplus_{p \in \Z} V(p)$ is a graded vector space and if the co-action $V \rightarrow \C[N] \otimes V$ is a graded homomorphism. If $N$ is connected, a graded $N$-module $V$ is equivalently defined as an $N$-module $V$ equipped with a $\Z$-grading such that, for any $i$ and $p$ in $\Z$, $\n_i \cdot V(p) \subseteq V(i + p)$.

By analogy, we call a \emph{filtered $N$-module} a filtered vector space $V$ equipped with a co-action $V \rightarrow \C[N] \otimes V$ that is a filtered homomorphism. If $N$ is connected, a filtered $N$-module $V$ is equivalently defined as an $N$-module $V$ equipped with a filtration such that, for any $i$ and $p$ in $\Z$, $\n_i \cdot V_p \subseteq V_{i + p}$. The filtration on $V$ induces a filtration on $V^N$ by setting $(V^N)_p \coloneq V_p \cap V^N$.

\begin{lemma} \label{lemma:gr-action}
	If $N$ is graded and $V$ is a filtered $N$-module, then $\gr V$ has an induced graded $N$-module structure.
\end{lemma}

\begin{proof}
	Taking the graded homomorphism associated with the filtered co-action gives the graded $N$-module structure on $\gr V$. 
\end{proof}

Assume that $N$ is connected. Recall that for an $N$-module $V$, the level $0$ of the Lie algebra cohomology computes the $N$-invariants:
$$ \Hgy^0(\n, V) = V^\n = V^N. $$
This cohomology is defined as the cohomology of the Chevalley--Eilenberg complex
\[ \Cpx^\bullet(\n, V) \coloneq V \otimes \operatorname{Alt}^\bullet \n^\vee. \] 
See \cite[Section 7.7]{weibel1994introduction} for precise definitions.

The following proposition is a generalisation of \cite[Proposition 5.2]{gan2002quantization} and \cite[Proposition 4.1.2.5]{juillard2025phd}. It gives a sufficient condition for the operations $\gr$ and $(\bullet)^N$ for a filtered module over a unipotent group $N$.

\begin{proposition} \label{proposition:gr-invariant}
	Let $N$ be a graded unipotent group whose Lie algebra is nonpositively graded: $\n = \bigoplus_{i \leqslant 0} \n_i$. Let $V$ be a filtered $N$-module such that:
	\begin{enumerate}
		\item the filtration on $V$ is bounded from below and exhaustive,
		
		\item there is an isomorphism $\C[N] \otimes W \cong \gr V$ of $N$-modules, where $\C[N]$ is equipped with the left-multiplication action and $W$ denotes a trivial $N$-module.
	\end{enumerate}
	Then there is a natural graded linear isomorphism 
	$$ \gr (V^N) \overset{\sim}{\longrightarrow} (\gr V)^N. $$
	Moreover, the Lie algebra cohomology $\Hgy^n(\n, V)$ is zero whenever $n \neq 0$. 
\end{proposition}

\begin{remark}
	In \cite{juillard2025phd}, we assumed that the Lie algebra grading was negatively supported: $\n = \bigoplus_{i \leqslant -1} \n_i$. One can relax this hypothesis to a nonpositively supported grading, but the unipotence of the group is no longer automatic.
\end{remark}

\begin{proof}
	Consider the dual $\Z$-grading on $\n^\vee$, that is nonnegative: $\n^\vee = \bigoplus_{i \geqslant 0} (\n^\vee)_i$. For $p$ in $\Z$ and $n \geqslant 0$, define the linear subspace
	$$ \Cpx(\n, V)^n_p \coloneqq \sum_{\substack{\mathcal{l}, i_1, \dots, i_n \in \Z \\ \mathcal{l} + i_1 + \cdots + i_n \leqslant p}} V_{\mathcal{l}} \otimes ((\n^\vee)_{i_1} \wedge \cdots\wedge (\n^\vee)_{i_n}). $$ 
	This defines a filtration on the complex such that the differential satisfies the inclusion $d(C^\bullet_p) \subseteq C^\bullet_p$ for all $p$ in $\Z$. According to \cite[Section~5.4]{weibel1994introduction}, there is a spectral sequence~$\{E_{r}\}_{r=0}^\infty$ associated with this filtered cochain complex.	
	
	For $p, q$ in $\Z$, the $0$-page is given by:
	$$ E_0^{p, q} = \gr_p C^{q-p} \coloneqq \bigoplus_{\mathcal{l} + i_1 + \cdots + i_{q-p} = p} \gr_\mathcal{l} V \otimes ((\n^\vee)_{i_1} \wedge \cdots\wedge(\n^\vee)_{i_{q-p}}). $$
	It is the Chevalley--Eilenberg complex for the $\n$-module $\gr V$: for $n \in \Z$, there is the identification
	$$ \bigoplus_{p \in \Z}  E_0^{p, n+p} = \gr V \otimes \operatorname{Alt}^n \n^\vee. $$
	Hence the $1$-page is the Lie algebra cohomology of $\gr V$:
	$$  \bigoplus_{p \in \Z}  E_1^{p, n+p} = \Hgy^{n}(\n, \gr V). $$ 
	
	By using the $\n$-equivariant isomorphism $\C[N] \otimes W \cong \gr V$, we get
	$$ \Hgy^{n}(\n, \gr V) \cong W \otimes \Hgy^{n}(\n, \C[N]). $$
	But the cohomology $\Hgy^{n}(\n, \C[N])$ coincides with the algebraic de Rham cohomology of $N$ as a consequence of the isomorphism between the cotangent sheaf of $N$ and $\C[N] \otimes \n^\vee$ \cite[Proposition 4.4.2]{springer2009linear}. Because the exponential map $\exp \colon \n \overset{\sim}{\rightarrow} N$ provides an isomorphism between $N$ and the affine space $\n$, its cohomology is given by
	$$ \Hgy^{n}(\n, \C[N]) \cong \delta_{n = 0} \, \C. $$
	Hence, the first page $E_1^{p, q}$ is zero unless $p=q$: the spectral sequence collapses.
	
	The filtration on $C$ is exhaustive and bounded from below because so is the filtration on $V$ and the dual of the Lie algebra has only nonnegative degree components. Hence, according to \cite[Theorem 5.5.1]{weibel1994introduction}, the spectral sequence converges to the Lie algebra cohomology of $V$: the infinity page is given by
	$$ E_\infty^{p, q} = \gr_p \Hgy^{q-p}(C^\bullet). $$	
	Because of the collapsing we get isomorphisms $E_\infty^{p, q} \cong E_1^{p, q}$. In particular, taking $p=q$, we see that the natural embedding
	$$  \gr (V^N) = \gr (V^\n) \hooklongrightarrow (\gr V)^\n = (\gr V)^N$$
	induced by the inclusion $V^N \subseteq V$ is an isomorphism. Taking $q = p + n$ for nonzero $n$, we get the cohomology vanishing.
\end{proof}

The following proposition states that $\gr$ and $(\bullet)^{\Gm}$ always commute for a filtered module over the multiplicative group.

\begin{proposition} \label{proposition:weight-filtration}
	Let $V$ be a filtered $\Gm$-module, where $\Gm$ is trivially graded. Then $\gr V$ is equipped with a natural $\Gm$-module structure that preserves the grading. 
	
	Denote by $V_p(\delta)$ and $(\gr_p V)(\delta)$ the weight spaces of the $\Gm$-actions, where $\delta$ and $p$ are in $\Z$. There is a linear isomorphism
	$$ \gr_p(V(\delta)) \cong (\gr_p V)(\delta) $$
	induced by the inclusion $V(\delta) \subseteq V$.
	
	For $\delta = 0$, this gives the isomorphism
	\[ \gr_p(V^\Gm) \cong (\gr_p V)^{\Gm}.  \]
\end{proposition}

\begin{proof}
	The first part of the statement is a consequence of \Cref{lemma:gr-action} applied to $\Gm$, seen as a trivially graded group so that $V$ is a filtered $\Gm$-module. 
	
	The filtration on $V(\delta)$ is defined by
	$$ V(\delta)_p \coloneqq V(\delta) \cap V_p = V_p(\delta). $$
	The inclusion $V(\delta) \subseteq V$ respects the filtrations so we get a graded inclusion~$\gr (V(\delta)) \subseteq (\gr V) (\delta)$. But $\gr_p V$ decomposes in two possible ways:
	$$ \bigoplus_{\delta \in \Z} \gr_p (V(\delta)) = \gr_p V = \bigoplus_{\delta \in \Z}  (\gr_p V) (\delta), $$
	hence the  inclusion $\gr_p (V(\delta)) \subseteq (\gr_p V) (\delta)$ is an equality. 
\end{proof}

